% !TeX root = ../Thesis.tex

%************************************************
\chapter{Diskussion}\label{ch:discussion} % $\mathbb{ZNR}$
%************************************************
\glsresetall % Resets all acronyms to not used

%************Beantwortung der Forschungsfragen**************
\section{Beantwortung der Forschungsfragen}
\label{sec:beantwortungfragen}
\begin{enumerate}
    \item \noindent Welche Grenzen zeigt ein Ansatz ohne münzspezifisches Training bei der Anwendung auf antike Münzbilder?

    \noindent Der Ansatz stößt an mehreren Stellen an Grenzen. Erstens liefert YOLO-World nur 23 (vorderseite) bzw. 22 (Rückseite) der 57 Prompts als Klasse mit meist niedriger Confidence (vgl. \cref{sec:ergebnisse}), und ob eine Klasse inhaltlich zutrifft, lässt sich ohne eigene Prüfung nicht sagen \todo{an qualitative Analyse anpassen}. Zweitens hängen Detektion und Vektoren von der festgelegten Prompt-Liste ab. Drittens [Einfluss von Erhaltungszustand und Bildtyp, aus Ergebnissen]. Viertens beschreibt ein Vektor das Merkmal ohne Kontext des Gesamtbildes. Zudem wird auf Münzen desselben Stempels häufig eine andere Klasse am sichersten detektiert, sodass nur 36\,\% der stempelgleichen Paare überhaupt verglichen werden (vgl. \cref{sec:ergebnisse}) Die Einzelheiten und die Grenzen der Untersuchung selbst behandelt \cref{sec:grenzen-untersuchung}.

\end{enumerate}

%*********Einordnung in bestehende Arbeiten**********************
\section{Einordnung in bestehende Arbeiten}
\label{sec:bestehendearbeiten}
Bildbasierte Verfahren für antike Münzen lassen sich nach ihrer Aufgabe in drei Gruppen einteilen. Die Arbeit liegt an der Schnittstelle dieser Gruppen, unterscheidet sich aber in einem Punkt von allen; sie verzichtet auf münzspezifisches Training.

\vspace{\baselineskip}

\noindent \textbf{Zuordnung ganzer Münzen} 

\noindent Guo et al. \cite{guo2023siamese} ordnen Münzbilder über ein Siamese-Netz mit Vision Transformern einer Prägung (Issue) zu, indem sie paarweise entscheiden, ob zwei Bilder zur selben Prägung gehören. Neue Prägungen lassen sich so ohne Neutraining ergänzen. Das Netz wird aber auf Münzpaaren trainiert, und verglichen wird die ganze Münze. Fare und Arandjelović \cite{fare2017ancient} untersuchen die Münzsuche systematisch im Hinblick auf den Erhaltungsgrad und nennen diesen eine zentrale Schwierigkeit. Der Erhaltungszustand spielt auch in dieser Arbeit eine Rolle (vgl. \cref{sec:qualitativeanalyse}).

\vspace{\baselineskip}

\noindent \textbf{Erkennung einzelner Motive und Zeichen} 

\noindent Cooper und Arandjelović \cite{cooper2020learning} trainieren je Motiv (z.\,B. Pferd, Adler, Schild) ein CNN, das erkennt, ob es auf der Rückseite vorkommt. Die Trainingslabels stammen aus Beschreibungstexten von Händlern, es wurden also keine Bildpositionen annotiert. Kim und Pavlovic \cite{kim2016discovering} suchen mit CNNs charakteristische Bildregionen je Münzklasse. Anwar \cite{anwar2026reading} behandelt einzelne Buchstaben als Objekte und trainiert verschiedene YOLO-Modelle nach, auf einem eigens annotierten Datensatz mit 38,808 Bounding Boxes. Das ist die engste Verbindung zu dieser Arbeit, denn auch dort wird ein YOLO-Detektor auf Münzbilder angewendet. Der Unterschied liegt darin, dass das Modell dort mit manuellen Annotationen nachtrainiert wird, hier jedoch ohne Nachtraining anhand von Textbeschreibungen detektiert. \todo{prüfen: 81\%, 84\%, NMI 0.959, 38.808 Boxen} \todo{anwar widerspruch: text 914\% und tabelle 90,4\%}

\vspace{\baselineskip}

\noindent \textbf{Stempelanalyse} 

\noindent Heinecke et al. \cite{heinecke2021unsupervised} gruppieren Münzbilder ohne Trainingsdaten nach dem Stempel, indem sie Schlüsselpunkte auf dem gesamten Münzbild vergleichen und die Gruppen statistisch bestimmen. Sie arbeiten auf Denaren aus der Zeit Neros und prüfen manuell, wie gut die Gruppen den echten Stempeln entsprechen. Das Ziel entspricht dem hier verwendeten Ansatz, Münzen desselben Stempels zu finden. Anders als hier werden aber keine einzelnen Merkmale detektiert und keine gelernten Feature-Vektoren verglichen.

\vspace{\baselineskip}

\noindent \textbf{Abgrenzung} 
\noindent Damit unterscheidet sich diese Arbeit in drei Punkten: Erstens kommt sie ohne Training und ohne Annotationen der Merkmalspositionen aus, wogegen die Verfahren zur Zuordnung und Motiverkennung \cite{guo2023siamese,cooper2020learning,anwar2026reading} auf münzspezifischen Daten trainiert werden. Zweitens wird nicht die ganze Münze verglichen, sondern einzelne lokalisierte Merkmale. Drittens dienen die Stempelgruppen aus dem Corpus Nummorum als Grundwahrheit für die Auswertung (vgl. \cref{sec:evaluationsaufbau}). Dafür wird auch weniger geleistet: Es werden nur 23 Klassen detektiert, die Stempelsuche wird auf 200 Münzen gezeigt, und eine inhaltliche Prüfung der Klassen fehlt \todo{an Ergebnisse anpassen}.

\vspace{\baselineskip}

\noindent \textbf{Vergleichbarkeit der Ergebnisse} 
\noindent Die Zahlen der genannten Arbeiten lassen sich nicht direkt mit den Ergebnissen dieser Arbeit vergleichen, da Aufgabe, Daten und Kennzahlen verschieden sind. Guo et al. berichten eine Genauigkeit von 81\,\% für die Zuordnung zu einer Prägung, Cooper und Arandjelović bis zu 84\,\% je Motiv, Anwar einen mAP50 von rund 90\,\% für Buchstaben. Heinecke et al. geben ein Clustermaß (NMI) von etwa 0,96 für Vorderseiten an. Diese Werte messen jeweils, ob etwas richtig eingeordnet wird. Hier wird gemessen, ob Münzen desselben Stempels unter den ähnlichsten Treffern liegen (HitRate@$k$, vgl. \cref{sec:evaluationsaufbau}). Die Ergebnisse dieser Arbeit liegen bei 17\,3\% (HitRate@1) und 46\,2\% (HitRate@5) \todo{aus \cref{sec:ergebnisse} einsetzen} und sind mit 200 Münzen als Tendenz zu verstehen. Ob ein Ansatz ohne Training den trainierten Verfahren nahekommt, lässt sich daraus nur eingeschränkt beantworten; ein Vergleich auf denselben Daten steht aus \todo{Formulierung nach den Ergebnissen anpassen}.

%***************Grenzen*****************************
\section{Grenzen}
\label{sec:grenzen}
\subsection{Grenzen des Ansatzes}
\label{sec:grenzen-ansatz}
\textbf{Bildmaterial} 

\noindent Die Bilder stammen teils von Originalen und teils von Gipsabgüssen und unterscheiden sich in Kontrast, Licht und Aufnahmebedingungen (vgl. \cref{sec:glnumismatik}). Zusammen mit Abnutzung und Korrosion kann dies die Ähnlichkeit zweier Münzen beeinflussen, unabhängig davon, ob sie vom selben Stempel stammen. Die Auswertung nach Bildtyp zeigt Unterschiede \todo{Ergebnis einsetzen}, trennt aber nicht von den Einflüssen des Erhaltungszustands.

\vspace{\baselineskip}

\noindent \textbf{Modell und Prompts} 

\noindent Die detektierten Klassen hängen von der Prompt-Liste mit 57 englischen Textbeschreibungen ab, die zu Beginn der Arbeit erstellt und nicht an die Motivbeschreibungen von CN angepasst wurde (vgl. \cref{sec:vorgehen}). Zudem hängen die Vektoren selbst von dieser Liste ab (vgl. \cref{sec:herausforderungen}), sodass ein anderes Vokabular zu anderen Ergebnissen führen kann. Ob eine detektierte Klasse inhaltlich zutrifft, wurde nur in der qualitativen Analyse an Stichproben geprüft (vgl. \cref{sec:qualitativeanalyse}). Der niedrige Schwellenwert von 0,005 führt zu vielen Detektionen mit geringer Confidence (vgl. \cref{sec:ergebnisse}). Auch der Grenzwert von 85\,\% der Bildfläche, ab dem Detektionen verworfen werden, wurde nicht systematisch bestimmt (vgl. \cref{sec:vorgehen}). 

Neben Yolov8s-World-v2 wurde Yolov8n verglichen, das auf der Vorderseite leicht bessere Trefferraten erzielt (vgl. \cref{tab:modelle_vorderseite,tab:modelle_rueckseite}). Die Unterschiede betragen fünf bis sieben von 200 Anfragen und sind daher gering. Der Modellvergleich beruht ausschließlich auf der Kosinus-Ähnlichkeit. Ob sich das Verhältnis der beiden Modelle mit der Manhattan-Distanz ändert, wurde nicht untersucht. \todo{falls doch, einfügen!}

\vspace{\baselineskip}

\noindent \textbf{Vektoren einzelner Merkmale} 

\noindent Der Vektor eines Merkmals wird aus dem ausgeschnittenen Bildausschnitt bestimmt und kennt den Kontext des Gesamtbildes nicht (vgl. \cref{sec:herausforderungen}). Ob Vektoren aus dem ganzen Bild mit Pooling auf der Feature-Karte bessere Ergebnisse liefern, wurde nicht untersucht.

\vspace{\baselineskip}

\noindent \textbf{Vergleich nur innerhalb einer Klasse} 

\noindent Merkmale werden nur mit Merkmalen derselben Klasse verglichen (vgl. \cref{sec:vorgehen}). Ob verwandte Klassen, die dasselbe Motiv unterschiedlich benennen, gemeinsam verglichen werden können, wurde im Rahmen dieser Arbeit nicht untersucht. Dadurch können Münzen desselben Stempels unverglichen bleiben, wenn auf ihnen verschiedene Klassen detektiert werden.

\subsection{Grenzen der Untersuchung}
\label{sec:grenzen-untersuchung}
\textbf{Umfang und Auswahl der Daten} 

\noindent Die Auswertung beruht auf 200 Münzen in 20 Stempelgruppen mit je zehn Münzen, die in aufsteigender ID-Reihenfolge und nicht zufällig aus dem Corpus Nummorum gewählt wurden (vgl. \cref{sec:gesamtansatzdatengrundlage}). Die Ergebnisse gelten daher für diese Auswahl und lassen sich nicht ohne Weiteres auf andere Prägungen, Epochen oder Regionen übertragen. Wegen der kleinen Gruppen könnten einzelne Münzen die Kennzahlen deutlich verändern.

\vspace{\baselineskip}

\noindent \textbf{Grundwahrheit} 

\noindent Als Grundwahrheit dienen die Stempelgruppen aus der CN-Datenbank \cite{cn-about}. Sie gilt nur für die Vorderseiten (vgl. \cref{sec:evaluationsaufbau}). Für die Rückseiten gibt es nur wenige Münzen mit gleichem Stempel, sie werden daher nur ergänzend ausgewertet. Die Güte der Auswertung hängt davon ab, dass die Zuordnung in CN korrekt ist. Sie wurde in dieser Arbeit nicht überprüft.

\vspace{\baselineskip}

\noindent \textbf{Kein Vergleich mit anderen Verfahren} 

\noindent Die Ergebnisse werden nur mit der Zufallsgrenze verglichen, nicht mit einem trainierten Modell oder einem klassischen Verfahren auf denselben Daten (vgl. \cref{sec:bestehendearbeiten}), außer mit YOLOv8n. Aussagen darüber, ob der Ansatz ohne Training mit trainierten Verfahren konkurrieren kann, sind daher nur eingeschränkt möglich.